\section{Agent Orchestration Frameworks}
In the landscape of modern artificial intelligence, an \textit{agent} is defined as an autonomous computational entity capable of sensing its environment, making informed decisions, and executing actions to achieve specific objectives. While standalone LLMs are powerful, their utility is often constrained by their inability to interact with the external world or maintain long-term state. To bridge this gap, orchestration frameworks are required to coordinate the interactions between models, external tools (e.g., compilers, hardware optimization engines), and persistent databases \cite{naveed2024comprehensiveoverviewlargelanguage, su2025daao}.

\subsection{LangChain: The Foundation of LLM Applications}
LangChain emerged as the first major framework to simplify the construction of LLM-powered applications. Its core philosophy centers on the concept of a "Chain"—a sequence of operations where the output of one step (e.g., a prompt) serves as the input for the next \cite{langchain}. By providing modular abstractions for prompts, memory, and tool integration (toolkits), LangChain enables developers to build complex pipelines. However, typical LangChain implementations are often linear or directed acyclic graphs (DAGs), which can limit the development of highly iterative workflows where an agent needs to loop back to a previous state based on evaluation results.

\subsection{LangGraph: Stateful and Iterative Orchestration}
To address the need for complex, cyclic workflows, LangGraph was developed as an extension of the LangChain ecosystem. Unlike traditional linear chains, LangGraph models agentic processes as a stateful graph where nodes represent computational steps and edges define the transition logic between them \cite{wang2024agentailanggraphmodular, hong2024workflowllm}. 

The primary innovation of LangGraph is its native support for \textbf{loops} and \textbf{state persistence}. This is particularly critical for MLOps tasks like "Neural Network Surgery," where an agent might need to modify a model, validate it against hardware constraints, and—if validation fails—revert to a previous architecture or adjust parameters. By maintaining a shared state that persists throughout the lifecycle of a request, LangGraph allows agents to maintain context across multiple turns and branch dynamically based on intermediate tool outputs.

\subsection{State Management and Fault Tolerance}
Reliable agentic systems require robust mechanisms to track history and recover from errors. State management in LangGraph is achieved through "checkpointing"—the practice of saving the entire state of the graph at each execution step. 
This enables several advanced features:
\begin{enumerate}
    \item \textbf{Human-in-the-loop (HITL):} Execution can be interrupted, allowing a human expert to review or modify the agent's plan before proceeding.
    \item \textbf{Time-travel Debugging:} Developers can inspect or resume execution from any previous state in the conversation history.
    \item \textbf{Automatic Recovery:} If a tool call or model invocation fails (e.g., due to rate limits or network issues), the system can resume from the last valid checkpoint without restarting the entire workflow \cite{yu2023stateful, wang2025stateful}.
\end{enumerate}
In this thesis, these capabilities are leveraged to ensure that the multi-stage conversion and optimization of STM32 models remain coherent and recoverable, even in the event of complex failures during the firmware generation or integration stages.

